\documentclass[11pt,a4paper]{article}
\usepackage[utf8]{inputenc}
\usepackage[T1]{fontenc}
\usepackage{lmodern}
\usepackage[french]{babel}
\usepackage{amsmath,amssymb,amsthm}
\usepackage[margin=2.2cm,top=2cm,bottom=2.2cm]{geometry}
\usepackage{enumitem}
\usepackage{fancyhdr}
\usepackage{array}
\usepackage[hidelinks]{hyperref}
\setlist[enumerate,1]{label=\arabic*.,leftmargin=*,itemsep=0.6em}
\setlist[enumerate,2]{label=(\alph*),leftmargin=1.8em,itemsep=0.3em}
\newcommand{\R}{\mathbb{R}}
\newcommand{\N}{\mathbb{N}}
\newcommand{\C}{\mathbb{C}}
\newcommand{\sh}{\operatorname{sh}}
\newcommand{\ch}{\operatorname{ch}}
\newcommand{\ee}{\mathrm{e}}
\newcommand{\dd}{\,\mathrm{d}}
\pagestyle{fancy}\fancyhf{}
\fancyhead[L]{Fiche 5 — Compléments}
\fancyhead[R]{sans indice, sans corrigé}
\fancyfoot[C]{\thepage}
\renewcommand{\headrulewidth}{0.3pt}
\newcommand{\exo}[3]{\vspace{1.1em}\noindent\textbf{\large #1}\hfill\textit{#2 \,·\, #3}\par\smallskip}
\begin{document}
\begin{center}
{\LARGE\bfseries Fiche 5 — Compléments}\\[0.3em]
{\large mardi 27 octobre · matin, 90 min}
\end{center}

\noindent\fbox{\parbox{\dimexpr\linewidth-2\fboxsep-2\fboxrule\relax}{\small Chronomètre lancé, sans cours ni notes, rédaction complète sur copie comme un oral ou un DS : hypothèses vérifiées, « à partir d'un certain rang », limites écrites avec leur flèche, aucun « $\sim$ » utilisé comme une égalité. Quand une question est finie, relis-la avec la \textbf{grille de rédaction}. Envoie-moi ta copie pour une correction ligne par ligne.}}


\exo{Exercice 1 — Équivalents de suites}{MPSI Camille Jullian 2023, ex. 1}{20 min}
\noindent Déterminer un équivalent simple de chacune des suites suivantes :
\begin{enumerate}[itemsep=0.4em]
\item $u_n=\dfrac{n^2+\ee^{-2n}+\sqrt{n^5}}{\ln(2n)+2n-3}$
\item $u_n=(n+3\ln(n))\ee^{-(n+1)}$
\item $u_n=\dfrac{\ln(n^2+1)}{n^2+1}$
\item $u_n=\ln\left(\dfrac{n^2+1}{n^2+2}\right)$
\item $u_n=\sqrt{n^2+n+1}-\sqrt{n^2-n+1}$
\item $u_n=\displaystyle\sum_{k=0}^{n}k!$
\item $u_n=\dfrac{n^{\sqrt{n+1}}}{(n+1)^{\sqrt{n}}}$
\end{enumerate}

\exo{Exercice 2}{MPSI Camille Jullian 2023, ex. 2}{20 min}
\noindent Soit $(u_n)$ une suite décroissante vérifiant $u_n+u_{n+1}\sim\dfrac1n$. Montrer que la suite converge nécessairement vers $0$ et en donner un équivalent simple. Le résultat reste-t-il vrai si la suite n'est pas supposée décroissante ?

\exo{Exercice 3}{MPSI Camille Jullian 2023, ex. 10}{20 min}
\noindent Montrer que les suites $(u_n)$ et $(v_n)$ définies par $u_n=n-\displaystyle\sum_{k=1}^{n}\cos\left(\dfrac1k\right)$ et $v_n=u_n+\sin\left(\dfrac1n\right)$ sont adjacentes (au moins à partir d'un certain rang).

\exo{Exercice 4}{MPSI Camille Jullian 2023, ex. 3}{30 min}
\noindent On considère la suite $(u_n)$ définie pour $n\geqslant 1$ par $u_n=\sqrt{n+\sqrt{n-1+\sqrt{n-2+\cdots+\sqrt{2+\sqrt{1}}}}}$.
\begin{enumerate}
\item Montrer que $(u_n)$ diverge vers $+\infty$.
\item Déterminer une relation simple entre $u_{n+1}$ et $u_n$.
\item Prouver par récurrence que $u_n\leqslant n$ puis que $u_n=o(n)$.
\item Déterminer un équivalent simple de $u_n$.
\item Déterminer $\displaystyle\lim_{n\to+\infty}u_n-\sqrt n$.
\end{enumerate}

\vfill
\begin{small}
\noindent\textbf{Sources.} MPSI : « Feuille d'exercices n$^{\circ}$17 : Analyse asymptotique », MPSI, lycée Camille Jullian, 14 mars 2023 (publiée sur normalesup.org/\textasciitilde glafon) ; les parenthèses d'indice des énoncés d'origine ont été retirées volontairement. Aucun énoncé n'a été réécrit.
\end{small}
\end{document}
